\chapter{Gradient Descent}
\labinfo{03\_Gradient\_Descent}{gradient\_descent.py}{UCI Concrete Strength: 1,030 mixes, 8 inputs}{Regression}{Stochastic GD test RMSE 6.656 MPa; $R^2=0.835$}

\section{Method}
Batch, mini-batch, and stochastic gradient descent were applied to a linear
strength model. Inputs were standardized from the training set; log age was
used to represent the age feature. The normal-equation solution provides a
reference optimum, and a learning-rate sweep demonstrates convergence behavior.

\section{Results}
Stochastic GD (learning rate 0.001) reached training loss 26.128 and test RMSE
6.656 MPa, close to the exact solution's 6.658 MPa. Batch GD at 0.1 was less
converged after 200 epochs (RMSE 6.839). A learning rate above the stability
limit diverged, illustrating the importance of step-size selection.

\begin{figure}[H]\centering
\includegraphics[width=.82\linewidth]{\figpath{03_Gradient_Descent}{gd_variants_loss.pdf}}
\caption{Training loss for batch, mini-batch, and stochastic updates.}
\end{figure}

\section{Python source}
\lstinputlisting[caption={Complete gradient descent implementation},label={lst:lab03}]{03_Gradient_Descent/gradient_descent.py}
